\section{Electronic Structure and Periodic Table}

\subsection{Light and Quantization}

\subsubsection{Photon energy}
\[
E=h\nu=\frac{hc}{\lambda}
\]
Photon energy is proportional to frequency and inversely proportional to wavelength. \(h\) is Planck's constant.

\subsubsection{Photoelectric effect}
\[
h\nu=E_{\mathrm{kin}}+W
\]
\(W\) is the work function of the metal. Electrons are emitted only when the photon energy exceeds the threshold.

\subsubsection{Bohr energy levels}
\[
E_n=-\frac{R_{\mathrm H}}{n^2}, \qquad n=1,2,3,\ldots
\]
Allowed electron energies in hydrogen are discrete. Transitions absorb or emit energy differences between levels.

\subsubsection{De Broglie wavelength}
\[
\lambda=\frac{h}{p}=\frac{h}{mv}
\]
Matter with momentum \(p\) has wave character. This motivates standing-wave orbitals for electrons.

\subsection{Quantum Numbers and Orbitals}

\subsubsection{Quantum numbers}
\[
\begin{array}{c c c}
\toprule
\text{symbol} & \text{allowed values} & \text{meaning}\\
\midrule
n & 1,2,3,\ldots & \text{shell}\\
\ell & 0,1,\ldots,n-1 & \text{subshell shape}\\
m_\ell & -\ell,\ldots,+\ell & \text{orbital orientation}\\
m_s & \pm \frac12 & \text{spin}\\
\bottomrule
\end{array}
\]
One electron in an atom is described by a unique set of four quantum numbers.

\subsubsection{Subshell labels}
\[
\ell=0,1,2,3 \quad \leftrightarrow \quad s,p,d,f
\]
For a subshell with angular momentum \(\ell\), there are \(2\ell+1\) orbitals.

\subsubsection{Orbital capacity}
\[
\text{maximum electrons in a subshell}=2(2\ell+1)
\]
Each orbital holds at most two electrons with opposite spin.

\subsection{Electron Configurations}

\subsubsection{Aufbau ordering}
\[
1s,\ 2s,\ 2p,\ 3s,\ 3p,\ 4s,\ 3d,\ 4p,\ 5s,\ 4d,\ldots
\]
Electrons occupy lower-energy orbitals first. The \(n+\ell\) rule gives the approximate filling order.

\subsubsection{Pauli exclusion principle}
\[
\text{no two electrons in an atom have the same }(n,\ell,m_\ell,m_s)
\]
The Pauli principle limits each orbital to two electrons.

\subsubsection{Hund's rule}
\[
\text{within degenerate orbitals, maximize unpaired parallel spins}
\]
Hund's rule gives the lowest-energy arrangement before pairing occurs.

\subsection{Periodic Trends}

\subsubsection{Effective nuclear charge}
\[
Z_{\mathrm{eff}}\approx Z-S
\]
\(Z_{\mathrm{eff}}\) is the nuclear attraction felt after shielding by other electrons. Across a period, \(Z_{\mathrm{eff}}\) generally increases.

\subsubsection{Atomic radius trend}
\[
\text{across a period: radius decreases}, \qquad
\text{down a group: radius increases}
\]
Increasing effective nuclear charge contracts atoms across a period. Additional shells increase radius down a group.

\subsubsection{Ionization energy}
\[
\ce{X(g) -> X+(g) + e-}, \qquad \mathrm{IE}>0
\]
Ionization energy is the endothermic energy required to remove an electron from a gaseous atom.

\subsubsection{Electron affinity}
\[
\ce{X(g) + e- -> X-(g)}
\]
Electron affinity describes the energy change when a gaseous atom accepts an electron.

\subsubsection{Mulliken electronegativity}
\[
\chi=\frac{\mathrm{IE}+\mathrm{EA}}{2}
\]
Electronegativity measures tendency to attract electron density in a bond. Fluorine is the most electronegative element in the lecture material.

\subsubsection{Isoelectronic ionic radius}
\[
\text{same electron count: larger }Z \Rightarrow \text{smaller radius}
\]
For isoelectronic species, increasing nuclear charge pulls the same electron cloud closer.
